\section{The five-branch boundary estimates}\label{sec:five-inputs}
Besides $f_n,\bar f_n,h_n$, write
\[
 W_n=A\bar f_{n-1}\quad(n\ge1),\qquad
 I_q(V)=\E[q^{-1/3}|V|^{8/3}].
\]
The arm message $W_n$ directs its first edge away from the root and
its remaining edges toward the first neighbor. Reversal exchanges
$f_n$ with $\bar f_n$ and $h_n$ with $W_n$.

\subsection{Simple-reference and ordinary moments}
The terminal residuals below contain at least two distinct lengths.
All messages include their actual equal-pair factors.
\begin{proposition}[Simple-reference fourth moments]\label{prop:five-reference}
For every existing outward arm $f_s$, $s\ge1$, every terminal arm
multiset as above, and every connector length $L\ge1$, there is a
structurally chosen terminal direction with actual received message
$X=T^LF$ such that
\[
 \E[f_sX^4]\le1.
\]
For $s\ge3$, the direction in the genuine-reference table satisfies
the stronger bound $1-d/3000$; for $s=2$ only the bound $1$ is used.
\end{proposition}
For $s=1$ the assertion is Theorem~\ref{thm:terminal} with step count
$L$. For $s\ge2$ the direction table and its proof are in
Sections~\ref{proof:R2-C032} and~\ref{proof:R2-C034}, indexed in
Appendix~\ref{app:five-directions}; the short-arm directions there are
specific to this reference catalogue.

For ordinary moments, a \emph{post block} means actual unequal central
pairs with negative canonical message
\[
 P=\prod_{i=1}^k W_{u_i}\bar f_{v_i},\qquad
 4\le u_1<v_1<\cdots<u_k<v_k.
\]
When nonempty, it is required that the designated anchor length satisfy
$m<u_1$. This order is automatic in the central allocations below.
\begin{proposition}[Ordinary eighth moments]\label{prop:five-ordinary}
Given an existing arm of length $m\ge3$, a terminal residual of size at
least two, a connector of length $L\ge1$ and a post block in the stated
order, there is a direction of exactly these resources with actual
message $H$ satisfying
\[
 \E H^8\le e^{-d/1024}.
\]
Given instead one existing unit arm and no unequal post block, there
is an actual message $H_1$ with $\E H_1^8\le1-d/20000$.
\end{proposition}
The complete ordinary direction table, including the first-edge
choices, is in Section~\ref{proof:R2-C046}, and
Appendix~\ref{app:five-directions} indexes its cases. The nonunit proof
draws on Sections~\ref{proof:R2-C030}, \ref{proof:R2-C038},
\ref{proof:R2-C039}, \ref{proof:R2-C040} and~\ref{proof:R2-C043}; the
unit input is the ordinary catalogue of Section~\ref{proof:Q5-C106}.
For the terminal $(1,2,3)$, the special direction occurs at $L=2$ in this
ordinary catalogue and at $L=1$ in the genuine-reference catalogue.

We also use the following path and mixed central bounds:
\begin{equation}\label{eq:five-mixed}
 \E f_s^2\le1\quad(s\ge2),\qquad
 \E[f_s^2Q^8]\le e^{-17d/64},
\end{equation}
where the second bound has
$Q=W_a\bar f_b\prod_iW_{u_i}\bar f_{v_i}$ and
$2\le s<a<b<u_1<v_1<\cdots$. The mixed proof, including $s=2$ and $a=3$,
is in Section~\ref{proof:R2-C044}. When $Q=1$ only the first bound is
needed. For $s=2$ its exact value is $1-d+\delta_2$; for $s\ge3$ the
finite path estimate gives $\E f_s^2\le1-d+2\delta_2\le1-3d/16$.

\subsection{A reference second moment for a compound branch}
A fork has a two-segment branch $c-r-t$ with independent lengths
$M,N\ge1$ and arbitrary side-arm tables at $r,t$.
\begin{proposition}\label{prop:five-compound}
For each $s\ge1$, this actual compound branch admits a direction with
received message $Y$ satisfying $\E[f_sY^2]\le1$.
\end{proposition}
\begin{proof}
If $t$ is reducible, its retained connector and possible residual arm
form one comparison arm at $r$. Choose the canonical endpoint direction
used in \eqref{q3:eq:singleendpoint}, apply that inequality, and restore
the equal pairs. Otherwise put
$G=T^{N-1}F^+_{t,N-1}$, so $\E[pG^4]\le1$.

If the residual at $r$ is empty or has at least two entries, select its
negative zero-step factor $C_r^-$ with $\E[q(C_r^-)^4]\le1$ and set
$H=C_r^-TG$, $Y=T^MH$. Row H\"older and Cauchy--Schwarz give
\begin{equation}\label{eq:inverse-half}
 (TG)^4\le q^3T(G^4),\qquad
 \E[H^2/q]\le\sqrt{\E[q(C_r^-)^4]\,
                         \E[(TG)^4/q^3]}\le1.
\end{equation}
Thus $H=\sqrt q\,h$ with $\|h\|_2\le1$. The operator bounds in
Proposition~\ref{q3:prop:operators}, together with normalized
propagation, give
\begin{equation}\label{eq:five-normalized}
 \sqrt{f_s}T^M\sqrt q
  =(\sqrt{f_s}T\sqrt q)(q^{-1/2}T\sqrt q)^{M-1},
 \qquad\|\sqrt{f_s}T^M\sqrt q\|\le1.
\end{equation}
This proves the claim in the R case, including $M=N=1$.

If $r$ has one residual arm of length $a$, orient it inward as
$\bar f_a$. At $t$, choose the positive canonical direction for an
even residual list and the negative canonical direction for an odd
residual list, and denote the resulting message by $S$. Let $S^-$
be its full reversal and put $V=A^NS^-$, $H=\bar f_aV$ and $Y=T^MH$.
The full reversal of \eqref{q3:eq:singleendpoint} gives
$\E[\bar f_aV^2]\le1$ for every $a,N\ge1$.
The cross-path contraction
$\|\sqrt{f_s}T^M\sqrt{\bar f_a}\|\le1$ proves the assertion.
Equal pairs at $r$ are restored by positivity.
\end{proof}

\subsection{The common-reference outlets for the path skeleton}
Consider one side $c-r-t$ of $P_5$. Its terminal connector $r-t$ has
length $L\ge1$, and its receiving connector $r-c$ has length $M\ge1$.
The terminal residual has size at least two, and the
residual at $r$ is nonempty. The following catalogue is used after the two
whole-tree alternatives of Section~\ref{sec:five-chain}.
\begin{proposition}[Common-reference outlet]\label{prop:five-outlet}
Except for a unit singleton at $r$ with $M=1$, every such side has an
actual message $U$ at $c$, determined by the structure, satisfying
\begin{equation}\label{eq:five-outlet}
 I_q(U)^3\le e^{-\eta d},\qquad \eta\ge1/500.
\end{equation}
Both sides of the path use the same $q$.
\end{proposition}
\begin{table}[htbp]\centering\small
\caption{The path outlet catalogue. The entries locate the proofs in the technical appendix, and Appendix~\ref{app:five-directions} indexes the direction tables.}\label{tab:five-outlets}
\begin{tabular}{@{}p{38mm}p{54mm}p{50mm}@{}}\toprule
Root residual & Structural subdivision & Where proved\\\midrule
Unit singleton & $M=2$; $M\ge3$ &
\ref{proof:R2-C026}, \ref{proof:R2-C029}, \ref{proof:R2-C031},
\ref{proof:R2-C033}, \ref{proof:R2-C035}, \ref{proof:R2-C037};
\ref{proof:Q5-C079}, \ref{proof:Q5-C106}\\
Singleton $a\ge2$ & $M=1,3,\ge4$; $a=2,M=2$; $a\ge3,M=2$ &
\ref{proof:R2-C032}, \ref{proof:R2-C034}, \ref{proof:R2-C045},
\ref{proof:Q5-E008-singleton}; \ref{proof:R2-C047}; \ref{proof:R2-C048}\\
Exactly two arms & Minimum one; minimum at least two &
\ref{proof:Q5-C066}, \ref{proof:Q5-C106}; \ref{proof:R2-C046}\\
Odd, at least three & Minimum one, two, or at least three &
\ref{proof:Q5-C185}, \ref{proof:Q5-C186}, \ref{proof:Q5-C188},
\ref{proof:Q5-C189}, \ref{proof:Q5-C192};
\ref{proof:R2-C002}--\ref{proof:R2-C005}\\
Even, at least four & Minimum one or at least two, with one actual inward multiplier retained &
\ref{proof:R2-C042}\\\bottomrule
\end{tabular}
\end{table}
The unit-singleton row $M=1$ is handled by the whole-tree construction
before this catalogue is consulted.
